When $\tilde{g}_u=\tilde{g}_d=\tilde{g}_s$ mixing between the $\tilde{\phi}$, the $\pi^0$ and the $\eta$ go to zero, causing the couplings between $\pi^0$ and $\eta$ and $\chi$ to go to zero, making these mesons unable to mediate interactions between the $\chi$ and the photon or decay into $\chi\bar{\chi}$.  This does not mean that the $\chi$ can no longer interact with the photon.  The $\tilde{\phi}$ can also mix with the $\eta'$.  We have generally neglected this meson because it is heavier than the others, and thus it's contribution to the cross section will be comparatively suppressed. We will consider it in this special case in which the coupling between the $\tilde{\phi}$ and the lighter neutral mesons goes to zero.   The $\eta$ and $\eta'$ mesons are actually the propagating mixed states of the $\eta_0$ and $\eta_8$ mesons.  $\eta_8$  appears in the meson octet, and mostly contributes to the $\eta$, $\eta_0$ is a singular ninth meson not part of the meson octet, that mostly contributes to the $\eta'$.  Throughout this work we have made the simplifying assumption that we can treat $\eta_8=\eta$ and $\eta_0=\eta'$ as each of these particles only gets a small contribution from the other particle it mixes with. Moving forward with these assumptions $\eta'$  appears in the chiral Lagrangian as $U = e^{\frac{i}{F}(\boldsymbol{\pi}+\eta')}$, expanding the chiral Lagrangian then gives the mixing term with the $\eta'$
The mixing term that appears in the Lagrangian is
\begin{equation}
    \mathcal{L}\supset-\frac{B_{0}F}{\sqrt{6}}(a_u+a_d+a_s)\eta'\tilde{\phi}\equiv -m_{\eta'\tilde{\phi}}\eta'\tilde{\phi},
\end{equation}

\begin{equation}
    \begin{split}
        \eta' =& \eta'_b \cos\theta-\tilde{\phi}_{b} \sin\theta\\ 
    \end{split}
\end{equation}\NAK{Is the \( \tilde{\phi} \) the bare one or the rotated one? I've assumed it to be bare given the context.}
This mixing results in the following interactions which allow the $\eta'$ to mediate interactions with $\chi$ and decay into $\chi$
\begin{equation}
\mathcal{L}\supset\tilde{g}_{\eta'\chi\chi}\eta'\chi\gamma_5\bar{\chi}+\tilde{g}_{\eta'\gamma\gamma}\eta'\epsilon^{\sigma\rho\mu\nu}\pi^0\partial_{\mu}A_{\sigma}\partial_{\nu}A_{\rho}
\end{equation}
% The fields have masses
% \begin{equation}
%     \begin{split}
%         m_1^2=&\frac{1}{2}\left(mP^2+m_{\eta'}^2+\sqrt{(mP^2-m_{\eta'}^2)^2+\frac{2}{3}B^2F_{\pi}^2(a_u+a_d+a_s)^2}\right)\\
%         m_2^2=&\frac{1}{2}\left(mP^2+m_{\eta'}^2-\sqrt{(mP^2-m_{\eta'}^2)^2+\frac{2}{3}B^2F_{\pi}^2(a_u+a_d+a_s)^2}\right)\\
%     \end{split}
% \end{equation}
The field couplings are
\begin{equation}
    \begin{split}
        \tilde{g}_{\eta'\chi\chi }=&-\tilde{g}_{\chi}\sin(\theta)\\
        \tilde{g}_{\tilde{\phi}\chi\chi }=&\tilde{g}_{\chi}\cos(\theta)\\
        \tilde{g}_{\eta'\gamma \gamma}= &\frac{e^2}{4\pi F_{\eta'}}\\
        \tilde{g}_{\tilde{\phi}\gamma \gamma}= &\tilde{g}_{\eta'\gamma \gamma}\tan(\theta)\\
        m_{\tilde{\phi}}^2&=\frac{m_{\tilde{\phi}_b}-m_{\eta}^2\sin(\theta)^2}{\cos(\theta)^2}
    \end{split}
\end{equation}
The mixing angle is
\begin{equation}
    \theta=\frac{1}{2}\text{ArcTan}\left(\frac{m_{\eta'\tilde{\phi}}}{m_{\phi_b}^2-m_{\eta'}^2+\frac{m_{\eta'\tilde{\phi}}}{4(m_{\phi_b}^2-m_{\eta'}^2)}}\right)
\end{equation}
In the limit that $m_{\eta'\tilde{\phi}}^2 \\ 4m_{\eta'}^2|(m_{\eta'}^2-m_{\tilde{\phi}}^2)|$
\begin{equation}
\begin{split}
    \tilde{g}_{\eta'\chi\chi }&\approx\frac{-\tilde{g}_{\chi}m_{\eta'\tilde{\phi}}}{2(m_{\phi_b}^2-m_{\eta'}^2)}   \\
    \tilde{g}_{\tilde{\phi}\chi\chi }&\approx\tilde{g}_{\chi}\\
    \tilde{g}_{\tilde{\phi}\gamma\gamma }&\approx\frac{\tilde{g}_{\eta'\gamma \gamma}m_{\eta'\tilde{\phi}}}{2(m_{\phi_b}^2-m_{\eta'}^2)}   \\
    m_{\tilde{\phi}}^2&\approx m_{\tilde{\phi}_b}^2\\
\end{split}
\end{equation},  This description breaks down if $m_{\eta'\tilde{\phi}}^2>4m_{\eta'}^2|(m_{\eta'}^2-m_{\tilde{\phi}}^2)|$
Here we approximate $F_{\eta'}$ using measurements of the $\eta'$ decay width to photons 
\begin{equation}
    F_{\eta'} = \sqrt{\frac{\alpha^2 m_{\eta'}^2}{64\pi^3\Gamma_{\eta'\rightarrow\gamma\gamma}}}
\end{equation}

These result in the following cross sections for scattering, annihlation to photons and production to photons respectively
\begin{equation}
    \frac{d\sigma_{\chi \gamma \rightarrow \chi \gamma}}{d\Omega}=\frac{- t^3}{256\pi^2 s }\left(\sum_j^{\eta', \tilde{\phi}}\frac{\tilde{g}_{j\chi\chi  } \tilde{g}_{j\gamma \gamma}}{(m_{j}^2-t)}\right)^2
    \end{equation}
    \begin{equation}
    \frac{d\sigma_{\chi \chi \rightarrow \gamma \gamma}}{d\Omega}=\frac{s^2}{256\pi^2 }\left(\sum_j^{\eta', \tilde{\phi}}\frac{\tilde{g}_{j\chi\chi  } \tilde{g}_{j\gamma \gamma}}{(m_{j}^2-s)}\right)^2\sqrt{\frac{s}{s-4m_\chi^2}}
\end{equation}
\begin{equation}
     \frac{d\sigma_{ \gamma \gamma\rightarrow \chi \chi}}{d\Omega}=\frac{s^2}{256\pi^2 }\left(\sum_j^{\eta', \tilde{\phi}}\frac{\tilde{g}_{j\chi\chi  } \tilde{g}_{j\gamma \gamma}}{(m_{j}^2-s)}\right)^2\sqrt{\frac{s-4m_\chi^2}{s}}
\end{equation}
This coupling will induce various new pion decays.  We present these below.
\begin{equation}
    \Gamma_{\eta'\rightarrow{\chi\bar{\chi}}} = \frac{\tilde{g}_{\eta'\chi\chi}^2\sqrt{\frac{m_{\eta'}^2}{4}-m_{\chi}^2}}{4\pi}
\end{equation}
\begin{equation}
    \Gamma_{k^+\rightarrow\pi^+\chi\bar{\chi}} = \frac{1}{32\pi^3m_{K^+}}\int_{m_{\pi^+}}^{E_{\pi \text{Max}}}dE_{\pi^+}(m_{K^+}^2+m_{\pi^+}^2-2m_{K^+}E_{\pi^+})Y(E_{\pi^+},m_{\chi})\left( \sum_j^{\eta', \tilde{\phi}} \frac{\tilde{g}_{j\chi\chi}\tilde{g}_{j\pi K}}{(m_{K^+}^2+m_{\pi^+}^2-2E_{\pi^+}m_{K^+}-m_{j}^2)}\right)^2
\end{equation}
where

\begin{equation}
\begin{split}
    \tilde{g}_{\tilde{\phi}\pi K}(E_{\pi^+})&=- (G_FF V_{us}V_{ud}^*)\sqrt{\frac{1}{3}}(m_{K^+}^2-m_{\pi^+}^2)(2g_8+3g_8^s)\sin(\theta)\\
     \tilde{g}_{\eta'\pi K} &= (G_FF V_{us}V_{ud}^*)\sqrt{\frac{1}{3}}(m_{K^+}^2-m_{\pi^+}^2)(2g_8+3g_8^s)\cos(\theta)
\end{split}
\end{equation}
\begin{equation}
    \Gamma_{K^+\rightarrow\pi^+\tilde{\phi}} = \frac{\tilde{g}^2_{\tilde{\phi}\pi K}}{32\pi m_{K^+}^3}\sqrt{(m_{K^+}^2+m_{\pi^+}^2-m_{\tilde{\phi}}^2)^2-4m_{K^+}^2m_{\pi^+}^2}
\end{equation}
For a scalar coupling to the $\chi$ these are
\begin{equation}
    \frac{d\sigma_{\chi \gamma \rightarrow \chi \gamma}}{d\Omega}=\frac{t^2(4m_{\chi}^2-t)}{256\pi^2 s }\left(\sum_j^{\eta', \tilde{\phi}}\frac{\tilde{g}_{j\chi\chi  } \tilde{g}_{j\gamma \gamma}}{(m_{j}^2-t)}\right)^2
    \end{equation}
    
    \begin{equation}
    \frac{d\sigma_{\chi \chi \rightarrow \gamma \gamma}}{d\Omega}=\frac{s(s-4m_{\chi}^2)}{256\pi^2 }\left(\sum_j^{\eta', \tilde{\phi}}\frac{\tilde{g}_{j\chi\chi  } \tilde{g}_{j\gamma \gamma}}{(m_{j}^2-s)}\right)^2\sqrt{\frac{s}{s-4m_\chi^2}}
\end{equation}
\begin{equation}
     \frac{d\sigma_{ \gamma \gamma\rightarrow \chi \chi}}{d\Omega}=\frac{s(s-4m_{\chi}^2)}{256\pi^2 }\left(\sum_j^{\eta', \tilde{\phi}}\frac{\tilde{g}_{j\chi\chi  } \tilde{g}_{j\gamma \gamma}}{(m_{j}^2-s)}\right)^2\sqrt{\frac{s-4m_\chi^2}{s}}
\end{equation}
This coupling will induce various new pion decays.  We present these below.
\begin{equation}
    \Gamma_{\eta'\rightarrow{\chi\bar{\chi}}} = \frac{\tilde{g}_{\eta'\chi\chi}^2(m_{\eta'}^2-4m_{\chi}^2)^{3/2}}{8\pi m_{\eta'}^2}
\end{equation}
\begin{equation}
    \Gamma_{k^+\rightarrow\pi^+\chi\bar{\chi}} = \frac{1}{32\pi^3m_{K^+}}\int_{m_{\pi^+}}^{E_{\pi^+ \text{Max}}}dE_{\pi^+}(m_{K^+}^2+m_{\pi^+}^2-2m_{K^+}E_{\pi^+}-4m_{\chi}^2)Y(E_{\pi^+},m_{\chi})\left( \sum_j^{\eta', \tilde{\phi}} \frac{\tilde{g}_{j\chi\chi}\tilde{g}_{j\pi K}}{(m_{K^+}^2+m_{\pi^+}^2-2E_{\pi^+}m_{K^+}-m_{j}^2)}\right)^2
\end{equation}
Kaon decay to $\tilde{\phi}+\pi^+$ is the same as for the all pseudo scalar coupling case.